% ---------------------------------------------------------------------------
%  06-limits.tex
% ---------------------------------------------------------------------------
\section{Limits}
\label{sec:limits}

\paragraph{It removes one exponential, not the exponential.}
\Cref{thm:decomp} attacks the product of \emph{independent} obligations. The
nesting towers of \cref{sec:problem} are not independent --- each level's
$\mathrm{AG}$ sits inside the previous level's $\mathrm{EF}$ and they share
atoms --- and their growth is untouched. Six separate attempts on that growth
are recorded in \cref{sec:negative}; none moved it. The remaining growth is
believed to be the genuine $2^{O(n)}$ of the bound.

\paragraph{The clock condition is a guard, not a heuristic.}
\ref{S2} refuses the split when a clock is constrained but owned by no
component, and it was worth asking whether a cleverer test --- dropping a clock
whose value no rule can still read --- would recover splits. It would recover
none. Counted over every family of both suites, that branch fires zero times
against $1590$ successful splits: dead-clock elimination has already freed
anything no pursuit and no deadline reads, so a clock still constrained at this
point is live and therefore owned. The condition stays as a guard against a
future rule that leaves a clock behind.

\paragraph{Splitting is not always worth it.}
The two families written so that decomposition cannot fire, \code{almost} and
\code{chain}, in fact split $436$ and $1099$ times: the shared atom stops
mattering once it has been consumed, and the residuals do fall apart. Both
families still measure below $1.0$ against the undecomposed search. So a split
that is available is not automatically a split that pays, and nothing here tries
to predict which is which.

\paragraph{The test can only see syntactic independence.}
Two components sharing a symbol are never split even when the symbol is used in
ways that do not interact. The symbols counted are atomic propositions and data
variables (\cref{sec:algorithm}); a semantic independence test would catch more
and cost more, and nothing here attempts one. Note which way the error runs: a
missed split costs time, an unsound one costs correctness, and the twelve
self-test failures on $\{x\le 5\}\wedge\{x\ge 10\}$ are what happens when the
symbol set is too small.

\paragraph{The suite is our own.}
The family that decomposes is one we wrote. See the caveat in
\cref{sec:eval}: the claim to make is about the mechanism, not about a speed-up
on a workload.

\paragraph{Witness extraction is disabled under decomposition.}
A witness for a decomposed node is the product of the components' automata, and
the extractor does not build it. With \code{--model} the static form declines to
decompose; the per-node form has not been reconciled with witness extraction at
all, which is the first thing to fix if either is turned on by default.

\paragraph{The overhead is not zero.}
$5$--$10\%$ on inputs that never decompose (\cref{sec:eval}). Turning either
form on by default needs either a cheaper negative test or evidence that real
inputs decompose often enough to pay for it.
